\section{Well-formedness}
\label{sec:well-formedness}

\PurePy excludes some Python behaviours statically and others only at runtime.
Well-formedness is the static part, a judgement on programs defined in this section; the operational
semantics (\secref{operational-semantics}) is defined for well-formed programs. A program is well-formed
when it is well typed, so the rules of this section give each expression a type, each statement a
static outcome recording either the assignments made or the fact that the statement definitely returns,
and each module a signature. The section introduces
types and subtyping, then follows the syntax of \secref{syntax}, taking statements, expressions and
patterns, then modules and programs, and \secref{operational-semantics} gives each of these a dynamic
counterpart.

\subsection{Key concepts}
\label{sec:principles}

This section introduces the concepts underlying the rules, each illustrated by admitted and rejected
programs; the rules follow in \secref{expression-typing} to \secref{well-formedness-modules}. The rules for
variables serve the two aims of \secref{no-mutable-state} and \secref{python-compatibility}, no mutable state and
compatibility with mypy.

\subsubsection{Variable declaration and definite assignment}
\label{sec:local-assignment}

\emph{Definite assignment} enforces the discipline of \secref{no-mutable-state}: a variable may be referenced only
where an assignment covers every path to the reference, and assigned only where no path has assigned it
already. Java imposes similar conditions on blank finals~\cite[\S16]{gosling2025}. The \ruleName{declare} rule binds the variable at $\DU{\tau}$, declared at $\tau$ and
definitely unassigned; the \ruleName{assign} rule requires that entry, checks the expression against $\tau$ and
makes the variable definitely assigned, so that $\Gamma(x)$ is a type, as the \ruleName{var} rule requires; an
annotated assignment does both at once; and the \ruleName{if} rules combine the branches with $\merge$, so
that a variable is definitely assigned after a conditional only if each branch either assigns it or definitely
returns, and is otherwise possibly unassigned, $\PU{\tau}$ (\secref{well-formedness-statements}). Branch-local variables are therefore permitted, but using one after the
conditional is not, nor is completing a partial assignment:

\begin{python}
def foo(b: bool) -> int:
    x: int
    if b:
        x = 3
        y: int = x + 1
        print(y)
    x = 0     # error: x assigned on the path through the branch
    return x
\end{python}

\noindent Python assigns \pyinline{x} twice when \pyinline{b} is true, so no pure reading of the program
agrees with Python: under shadowing, \pyinline{return x} would refer to a different variable from the one
assigned in the branch.

A function defined in a branch can also be used only within that branch; its name is possibly unassigned after the
conditional, and no other branch can complete the definition, because a definition is also a declaration and a scope
declares a name at most once. A variable of \pyinline{Callable} type, declared before the conditional and assigned in
each branch, can be used if the function is needed after the conditional.

The \ruleName{def-body} rule introduces every variable declared or assigned by the body, $\assignsS(s)$, at $\Unbound$
before checking the body, so a reference before the declaration is rejected statically, whereas Python fails at
run time with \pyinline{UnboundLocalError}:

\begin{python}
y: int = 7
def f() -> int:
    x: int = y  # error: y is local, because of the declaration below
    y: int = 8
    return x
\end{python}

\noindent Here $\assignsS(s)$ is the set of variables declared or assigned anywhere in a statement $s$, not
descending into nested function definitions, which have their own scope.

The \ruleName{module} rule does the same for a module body, with the builtins as the enclosing scope
(\secref{scope}), so a use of a builtin name is rejected if the module later declares that name:

\begin{python}
n: int = len([1])  # error: len is bound by the module
len: int = 2
\end{python}

\subsubsection{Variable capture}
\label{sec:captured-variables}

In Python a closure holds a mutable binding, so a variable captured by a lambda or nested function
definition may be reassigned afterwards, changing the value seen by the function (\secref{no-mutable-state}).
This cannot arise in \PurePy: a closure captures only definitely assigned variables, and a definitely
assigned variable is never assigned again. One case remains, in comprehensions: in Python a comprehension mutates its generator
variables on each iteration, so \PurePy forbids a lambda in the body or in a later qualifier from capturing a generator
variable, a side condition of \ruleName{list-comp}, \ruleName{dict-comp} and \ruleName{qual-generator}.

Here $\captures(\seq{g})$ is the set of variables of the enclosing scope captured by the lambdas within the
qualifiers $\seq{g}$ of a comprehension, with $\captures(e)$ computing the same for an expression $e$: the
free variables of each lambda's body, less its parameters.

\subsubsection{Mutual def regions}
\label{sec:mutual-recursion}

Mutual recursion in Python relies on late binding: the free variables of a function are resolved when
they are evaluated, so two functions may refer to each other whatever order they are defined in, provided
both exist by the time the reference is evaluated. \PurePy makes the
dependency static: a \emph{mutual def region}, a maximal contiguous block of \pyinline{def} statements (\figref{syntax-stmt}), is a single simultaneous binding. The \ruleName{defs} rule checks each body in a
context in which every name of the mutual def region is definitely assigned, so any function in the region may refer
to any other, and requires the names to be distinct. Self-recursion is the special case of a singleton
mutual def region. Any other statement between two definitions separates them into different mutual def regions:

\begin{python}
def even(n: int) -> bool:
    if n == 0:
        return True
    return odd(n - 1)  # error: odd is not yet bound

x: bool = even(0)  # intervening statement separates the mutual def regions

def odd(n: int) -> bool:
    if n == 0:
        return False
    return even(n - 1)
\end{python}

\noindent This is rejected statically: \pyinline{even} forms a mutual def region by itself, in which \pyinline{odd}
is not yet bound, so the reference fails the \ruleName{var} rule. Python accepts the program, since the
call \pyinline{even(0)} never reaches \pyinline{odd}. A mutual def region with two definitions of the same name is
rejected by the distinct-names condition:

\begin{python}
def g() -> int:
    return 0

def f() -> int:
    return g()
def g() -> int:  # error: g defined twice in the same mutual def region
    return 1
\end{python}

\noindent Python accepts this too, treating the second \pyinline{def g} as a reassignment of \pyinline{g},
which late binding then makes the target of the call in \pyinline{f}. Finally, the branches of a conditional may contain mutual def regions, with any use of those
functions after the conditional subject to the usual definite assignment rules (\secref{local-assignment}).

\subsubsection{Pattern binding and reachability}
\label{sec:pattern-binding}

A \pyinline{case} pattern binds variables, and Python scopes them as it does assignments: they remain in
scope after the \pyinline{match} statement. \PurePy follows: each case is treated as a branch, and its pattern as an
assignment of the variables it binds, which must not be declared already. The set of variables a pattern $p$ introduces is $\binds(p)$; the
\ruleName{match} rule checks each body in a context in which they are definitely assigned, and combines
the cases with $\merge$ as the \ruleName{if} rules do, so a variable is definitely assigned after the match
only if every case binds it or returns. A match whose cases leave some value of the scrutinee type
unmatched may fall through; in this case the \ruleName{match-partial} rule merges in the empty assignment
set, as for a conditional without an \pyinline{else}:

\begin{python}
v: tuple[int, int] = (1, 2)
match v:
    case (x, y):
        print(x, y)
print(x)  # error: x not definitely assigned
\end{python}

\noindent Python prints \pyinline{1}, the binding having escaped the match. Within a pattern the
variables must be distinct, as in Python. Every case must also be reachable: the matching judgement rejects a case
that can match only values already matched by an earlier case (\secref{well-formedness-patterns}), which
Python leaves silently unreachable:

\begin{python}
v: tuple[int, int] = (1, 2)
match v:
    case (a, b):
        print("seq", a, b)
    case (1, 2):  # error: the previous case matches every pair
        print("lit")
\end{python}

\subsubsection{Reachability}
\label{sec:unreachable}

Python accepts statements that can never run, such as code after a \pyinline{return}. \PurePy rejects
them, the statement-level counterpart of case reachability (\secref{pattern-binding}). A statement
sequence returns only if its last statement returns and every earlier one assigns: the \ruleName{seq}
rule requires the first statement of a sequence to have static outcome $\Assigns{\Delta}$
(\secref{well-formedness-statements}), so a statement following a \pyinline{return} has no derivation:

\begin{python}
def foo(x: int) -> int:
    return x + 1
    return x  # error: unreachable
\end{python}

\subsubsection{Module checking}
\label{sec:load-order}

In Python a module's namespace is populated as loading proceeds, by the assignments in its body and, for
a module with submodules, by the import system binding each submodule as it loads, so the referent of a qualified name can
depend on which modules have loaded so far, anywhere in the program (\secref{no-load-dependent}). \PurePy gives each name a meaning fixed by the
text and imports of the module alone. Imports come first in a module, so its dependencies are known before any
statement runs. Loading is checked statically: the \ruleName{module} rule checks a body and computes its
signature, and the \ruleName{import} and \ruleName{from-import} rules check each target and its
containing modules first, so an import cycle has no derivation. A module may not bind a name that is also the name of one of its
submodules, a side condition of \ruleName{module}, since Python resolves the clash by load order. A
qualified reference must resolve through the imports of the referring module to a definitely assigned
member. A submodule is not an attribute of its parent until it is loaded, so
\pyinline{import pkg} followed by \pyinline{pkg.sub.y} is rejected unless the module also imports
\pyinline{pkg.sub}, whereas Python permits the reference whenever some earlier import, anywhere in the
program, happens to have loaded \pyinline{pkg.sub}:

\begin{minipage}[t]{0.3\textwidth}
\begin{python}
# a
import pkg.sub
\end{python}
\end{minipage}
\hfill
\begin{minipage}[t]{0.6\textwidth}
\begin{python}
# __main__
import pkg
import a
print(pkg.sub.y)  # error: pkg.sub is not imported here
\end{python}
\end{minipage}

\noindent Python accepts this because importing \pyinline{a} loads \pyinline{pkg.sub}; remove that import,
which the main module does not otherwise use, and Python raises \pyinline{AttributeError}. It raises the
same for a reference to a module member never assigned by the body, such as \pyinline{lib.x} where
\pyinline{lib} assigns \pyinline{x} only on a branch that is not taken. \PurePy rejects all three programs
statically (\secref{well-formedness-modules}).



\subsection{Types}
\label{sec:types}

Well-formedness is checked relative to a context $\Gamma$, recording what each name in scope denotes, and a
class table $\Sigma$, recording the classes defined so far (\secref{contexts}). Programs mention types only
through the type expressions in annotations (\secref{types-schemes}), so each typing rule must resolve these to
types by looking up the names in $\Gamma$. Class names applied to type arguments resolve to
instantiated classes, alias names to the bodies of the aliases at their type arguments, and type parameters to type
variables; the full definition is given in \specfigref{resolve-type}.

\paragraph{Subtyping.}

Subtyping (\specfigref{subtyping}) is structural and follows mypy (\secref{python-compatibility}). Instantiated classes are below their base classes, with type arguments invariant, as
are lists and dictionaries; tuples are covariant, and callables are contravariant in their parameters and covariant
in their results. A literal type is below its base type, and \pyinline{int} is below \pyinline{float}, as PEP~484
requires, whereas \pyinline{bool} is not below \pyinline{int}, although in Python \pyinline{bool} is a subclass of \pyinline{int}. \replaced{A union $\tyUnion{\sigma}{\tau}$ is a least upper bound of its \emph{disjuncts} $\sigma$ and $\tau$}{A union is above
each of its disjuncts and below every type above all of them}; \pyinline{Never} is the least type and
\pyinline{object} the greatest. Subtyping is a preorder (Lemma~\specnum{lem:subtyping-transitive}), with distinct
types such as $\tyUnion{\sigma}{\tau}$ and $\tyUnion{\tau}{\sigma}$ \emph{equivalent}, $\sigma \tyEquiv \tau$, when each
is a subtype of the other; the notations $\subty$ and $\tyEquiv$ follow \citet{pierce2002}.

\begin{specdefinition}{Comparability, join and meet}
Two types are \emph{comparable}
when either is a subtype of the other. Two comparable types
\emph{join} to the larger type and two incomparable types to their union. Two comparable types \emph{meet} to the smaller type; a
meet of incomparable types distributes over a union, takes tuples componentwise and callables with
parameters joined and results combined using meet, and is \pyinline{Never} otherwise, since no value belongs to both types.
So joining a type with itself gives that type, joining a class with its base class gives the base class, and
joining \pyinline{int} with \pyinline{str} gives $\tyUnion{\tyInt}{\tyStr}$. A sequence of types is joined by
folding binary join over the sequence, with \pyinline{Never} as the unit\speconly{ (\S\ref{aux:join})}.
\end{specdefinition}

\begin{lemma}[Type lattice]
\label{lem:lattice}
Types ordered by subtyping in $\Sigma$, up to equivalence, form a bounded distributive lattice: $\joinTy{\Sigma}{\sigma}{\tau}$
is the least upper bound of $\sigma$ and $\tau$, $\meetTy{\Sigma}{\sigma}{\tau}$ is their greatest lower bound,
\pyinline{object} is the greatest type and \pyinline{Never} the least.
\end{lemma}


\subsection{Expressions}
\label{sec:expression-typing}

\replaced{Typing of expressions is bidirectional, with two judgements.}{Within an expression the typing judgement gives each subexpression a type, and
is bidirectional.} \emph{Synthesis} $\Sigma; \Gamma \vdash e \synth \tau$ reads a type out of an expression (\specfigrefThree{typing-synth}{typing-synth-elim}{typing-synth-ctd}),
and \emph{checking} $\Sigma; \Gamma \vdash e \checks \tau$ tests an expression against a type supplied by its
surroundings (\specfigref{typing-check}). The division follows the usual recipe for bidirectional
systems~\citep{davies2000,dunfield2021}\replaced{, in which introduction forms check and elimination forms synthesise, since an elimination form starts from a subterm whose type is known.}{: introduction forms check, since a lambda's parameters cannot be
annotated and an empty list does not determine its element type, and elimination forms synthesise, since a call, a subscript or an attribute
reference starts from a subterm whose type is known.} \replaced{Such systems usually provide an annotation form $(e \mathbin{:} \tau)$, which lets a checking form occur
where synthesis is required. Python does not have an annotation form, so}{Python has no annotation form
$(e \mathbin{:} \tau)$, by which such systems usually let a checking form occur where synthesis is required, so}
\PurePy gives each introduction form a synthesis rule wherever its parts determine the type, an approach also taken
by Swift, Kotlin and Scala~\citep{swift2025,kotlin2025,odersky2025}. \added{Lambdas and empty lists have only checking rules, since Python does not allow lambda
parameters to be annotated, and empty lists fail to determine their element types. Lambdas take their parameter types
from the expected \pyinline{Callable} type, and check their bodies against the return type:
\begin{mathpar}
\small
\inferrule*[lab={\ruleName{lambda}}]
{
  \Sigma; \Gamma \override \set{\bind{\seq{x}}{\seq{\tau}}} \vdash e \checks \tau
}
{
  \Sigma; \Gamma \vdash \exLambda{\seq{x}}{e} \checks \tyCallable{\seq{\tau}}{\tau}
}

\end{mathpar}
}\added{Literals other
than floats synthesise their literal types, such as $\tyLiteral{\pyinline{3}}$, and lists synthesise the join of
the base types of their elements, since \pyinline{list} is invariant and a literal element type would rule out every other
element. Expressions that synthesise $\sigma$ check against any supertype of $\sigma$ by subsumption, which applies only
when no other rule concludes $\Sigma; \Gamma \vdash e \checks \tau$, so both judgements are syntax-directed and
expressions have at most one synthesis derivation, and at most one checking derivation against a given type.}

\begin{DeletedBlock}
Variables, literals, attribute references, calls, constructor calls, operators and subscripts synthesise.
A variable synthesises a type only when its context entry is a type, so a variable that is not definitely
assigned, or a name bound at a type scheme, does not synthesise a type. A
literal $\ell$ synthesises the literal type $\tyLiteral{\ell}$, except for a float literal, which synthesises \pyinline{float}.

In a subscript $\exSubscript{e}{e'}$ there is a rule for each container that $e$ can synthesise: a list or string with $e'$ of type \pyinline{int}, and a
dictionary with $e'$ of type \pyinline{str}. Tuples have three rules: an index of type \pyinline{int} gives
the join of the component types, unless it synthesises a literal type, in which case it
gives the component type at that position, counting from the end when the index is negative:
\begin{mathpar}
\small
\inferrule*[lab={\ruleName{subscript-tuple}}]
{
  \Sigma; \Gamma \vdash e \synth \tyTuple{\tau_1, \ldots, \tau_m} \\
  \Sigma; \Gamma \vdash e' \synth \tyLiteral{n} \\
  0 \leq n < m
}
{
  \Sigma; \Gamma \vdash \exSubscript{e}{e'} \synth \tau_{n+1}
}

\end{mathpar}
\end{DeletedBlock}
\added{Operators are overloaded, as specified in \specfigref{operator-overloads}, and the \ruleName{binop} and
\ruleName{unop} rules synthesise the return type of the most specific overload at the operand types. }\replaced{Union types need dedicated rules at elimination forms, operators included: with \pyinline{x} of type
$\tyUnion{\tyInt}{\tyStr}$, no overload of \pyinline{*} applies at $\exBinOp{\pyinline{x}}{\pyinline{*}}{\pyinline{2}}$, which is
instead typed at \pyinline{int} and at \pyinline{str} and synthesises $\tyUnion{\tyInt}{\tyStr}$. Calls in which the callee
synthesises a union are typed at each disjunct, with the callee replaced by a fresh variable bound at that disjunct,
and synthesise the join of the two results:}{A subscript, attribute reference or call with a subject synthesising a union is typed at each disjunct, with a variable in place of the subject, and synthesises the
join of the results:}
\begin{mathpar}
\small
\inferrule*[lab={\ruleName{call-union}}]
{
  \Sigma; \Gamma \vdash e \synth \tyUnion{\sigma}{\sigma'} \\
  \Sigma; \Gamma \override \set{\bind{y}{\sigma}} \vdash \exCall{y}{\seq{e}} \synth \tau \\
  \Sigma; \Gamma \override \set{\bind{y}{\sigma'}} \vdash \exCall{y}{\seq{e}} \synth \tau' \\
  y \notin \fv(\seq{e})
}
{
  \Sigma; \Gamma \vdash \exCall{e}{\seq{e}} \synth \joinTy{\Sigma}{\tau}{\tau'}
}

\end{mathpar}
\deleted{Generators need no such rule, because $\elemTy{\Sigma}$ takes a union to the join of its disjuncts' element types.
}\replaced{Subscripts and attribute references have analogous rules. A single rule for an arbitrary evaluation
context~\citep{dunfield2004} can instead eliminate a union at any subterm in evaluation position, including an argument,
which has the advantage of accepting programs in which a function of intersection type is applied to an argument of union
type, each disjunct needing a different conjunct of the function type. However, the rule does not determine where it applies, so
type checking requires search~\citep{dunfield2011}. \PurePy lacks intersection types, and the rules above cover
overloaded operators, which are their closest analogue; stating union elimination separately for
each elimination form therefore keeps both judgements syntax-directed without rejecting any program which that rule
would accept.}{\citet{dunfield2004} state union elimination once, for an arbitrary evaluation context; here union elimination is
stated separately for each elimination form. }

\begin{DeletedBlock}
The \ruleName{call} rule for a call $\exCall{e}{\seq{e}}$ uses the auxiliary
judgement $\Sigma; \Gamma \vdash e, \seq{e} \synth \seq{\sigma}, \sigma$ (\specfigref{typing-callee}), which determines parameter
types $\seq{\sigma}$ and return type $\sigma$ from the callee along with its arguments. The remaining premises check
the arguments against $\seq{\sigma}$, and the call synthesises $\sigma$. The auxiliary judgement has two rules. A callee that synthesises a
\pyinline{Callable} type gives the parameter and return types of that type.
\end{DeletedBlock}
\added{Generic functions, bound in the context at type schemes rather than types, do not have types on their own,
but are typed only at calls or against expected types. Calls of generic functions }\deleted{A name bound at a
type scheme $\scheme{\seq{\alpha}}{\tyCallable{\seq{\sigma}}{\sigma}}$ }\replaced{instantiate}{instantiates} the type parameters from the arguments\deleted{
(\ruleName{callee-scheme})}. The type arguments are \replaced{determined by matching each parameter type against the type synthesised by the corresponding argument}{collected by $\typeArgs{\Sigma}$ over the arguments that synthesise,
by a one-way unification of each parameter type with the synthesised type, in which the type parameters $\seq{\alpha}$ are the
unknowns: a type parameter takes the type found at the same
position in the synthesised type}, with \replaced{a type parameter}{one} found at several positions taking the join of those types\added{,
and every argument is then checked against its instantiated parameter type:}\deleted{.
Arguments that do not synthesise, such as lambdas, are left out. The parameter and return types are
instantiated at the type arguments given by the context, which must bind every type parameter:}
\begin{DeletedBlock}
\begin{mathpar}
\small
\inferrule*[lab={\ruleName{callee-scheme}}]
{
  \Gamma \vdash q : \scheme{\seq{\alpha}}{\tyCallable{\seq{\sigma}}{\sigma}} \\
  I = \set{i \mid \Sigma; \Gamma \vdash e_i \synth \tau_i} \\
  \typeArgs{\Sigma}(\seq{\alpha}, (\sigma_i \mid i \in I), (\tau_i \mid i \in I))(\seq{\alpha}) = \seq{\tau}'
}
{
  \Sigma; \Gamma \vdash q, \seq{e} \synth \subst{\seq{\tau}'}{\seq{\alpha}}{\seq{\sigma}}, \subst{\seq{\tau}'}{\seq{\alpha}}{\sigma}
}

\end{mathpar}
\end{DeletedBlock}
\begin{AddedBlock}
\begin{python}
def keep[T](p: Callable[[T], bool], xs: list[T]) -> list[T]:
    return [x for x in xs if p(x)]
ys: list[int] = keep(lambda x: x > 0, [1, -2, 3])
\end{python}

\noindent The call fixes \pyinline{T} at \pyinline{int} from the list, and the lambda is then checked against
$\tyCallable{\tyInt}{\tyBool}$. Checking an argument never yields type arguments, which lets all type arguments
be determined in one step, before any argument is checked, so a call is rejected if a type parameter occurs only in the types of parameters with arguments that do not
synthesise, such as lambdas. These calls are accepted by mypy, which types each lambda once the other type arguments are known; the
Python typing specification leaves the inference of type arguments in calls of generic functions entirely
unspecified. Generic functions check against a type $\tau$ by
instantiating their type parameters from $\tau$ in the same way, provided the instantiated type is a subtype of $\tau$,
so they can be passed where \pyinline{Callable} types are expected.
\end{AddedBlock}

\deleted{Lambdas, applied lambdas, lists, tuples, dictionaries, conditional expressions and comprehensions check.
}\deleted{A lambda takes its parameter types from the expected \pyinline{Callable} type, and its body checks against the return type:}
\deleted{A name bound at a type scheme $\scheme{\seq{\alpha}}{\sigma}$ checks against a type $\tau$
by instantiating $\seq{\alpha}$ at the type arguments collected by $\typeArgs{\Sigma}$ from $\sigma$ and $\tau$, provided the
instantiated $\sigma$ is below $\tau$, so a function with type parameters can be passed where a \pyinline{Callable} type is expected.
An expression that synthesises $\sigma$ checks against any supertype
of $\sigma$, which applies only when no other rule concludes $\Sigma; \Gamma \vdash e \checks
\tau$, so both judgements are
syntax-directed and an expression has at most one derivation in each.}

\begin{DeletedBlock}
\paragraph{Operators.}

Each operator is overloaded as specified in \specfigref{operator-overloads}.

The \ruleName{binop} and \ruleName{unop} rules (\specfigref{typing-synth}) take the return type
of the most specific overload at the operand types, $\resolveOp{\Sigma}$, so \pyinline{1 + 2} is an
\pyinline{int} and \pyinline{1 + 2.0} a \pyinline{float}.
\end{DeletedBlock}

\subsection{Statements}
\label{sec:well-formedness-statements}

\begin{specdefinition}{Static outcome}
A well-typed statement has a \emph{static outcome} $R$, which is either $\Returns$ or $\Assigns{\Delta}$. The
outcome is $\Returns$ when the statement \emph{definitely returns} or definitely reaches \pyinline{assert False}.
Otherwise it is $\Assigns{\Delta}$, where $\Delta$ is a context with an entry for each variable declared or assigned
by the statement (Definition~\ref{def:context}), recording whether the variable is definitely assigned, definitely
unassigned or possibly unassigned after the statement.
\end{specdefinition}

Context operations $\override$ and $\merge_{\Sigma}$ lift to a right semiring over static outcomes $R$, with $\Assigns{\varnothing}$ as left and right unit for $\override$ and $\Returns$ as unit for $\merge_{\Sigma}$ and zero for $\override$.

\begin{definition}[Override and merge on static outcomes]
\begin{defalign}
\Assigns{\Delta} \merge_{\Sigma} \Assigns{\Delta'} & = \Assigns{\Delta \merge_{\Sigma} \Delta'} \\
R \merge_{\Sigma} \Returns = \Returns \merge_{\Sigma} R & = R \\
\\
\Assigns{\Delta} \override \Assigns{\Delta'} & = \Assigns{\Delta \override \Delta'} \\
R \override \Returns = \Returns \override R & = \Returns
\end{defalign}
\end{definition}

\noindent The $\Returns \override R$ case never arises because statement typing excludes unreachable statements\paperonly{ (\secref{unreachable})}.


The statement judgement $\Sigma; \Gamma; \zeta \vdash s : R$ says that under $\Gamma$ and
optional return type $\zeta$, statement $s$ has static outcome $R$ (\specfigref{typing-stmt}); every rule passes $\zeta$ down unchanged.
A sequence $\exSeq{s}{s'}$ types $s$, then types $s'$ under $\Gamma$ overridden by the context assigned
by $s$. A conditional combines its
branches with $\merge_{\Sigma}$, merging in $\Assigns{\varnothing}$ when there is no \pyinline{else} branch.
A call statement with a callee synthesising a union is typed
once per disjunct (\ruleName{call-stmt-union}), following \ruleName{call-union} for a call expression.

A \pyinline{match} scrutinee synthesises a type $\sigma$, and its cases $\seq{b}$ are checked in turn by the
auxiliary judgement $\Sigma; \Gamma; \zeta \vdash \tau, L, \seq{b} : S, \tau', L'$ (\specfigref{typing-cases}), which threads
a remaining type, starting from $\sigma$, and a residual $L$ of shapes, starting from the shapes of $\sigma$ (\figref{shapes}),
through the cases, merging their outcomes into $S$. The \ruleName{match} rule applies when the final remaining type is $\tyNever$, meaning the match is
exhaustive; the \ruleName{match-partial} rule applies otherwise and also merges in $\Assigns{\varnothing}$, the
outcome of the path on which the scrutinee is not matched by any of the cases:
\begin{mathpar}
\small
\inferrule*[lab={\ruleName{match}}]
{
  \Sigma; \Gamma \vdash e \synth \sigma \\
  \Sigma; \Gamma; \zeta \vdash \sigma, \shapes{\Sigma}(\sigma, \varnothing), \seq{b} : S, \tyNever, \varnothing
}
{
  \Sigma; \Gamma; \zeta \vdash \exMatch{e}{\seq{b}} : S
}

\and
\inferrule*[lab={\ruleName{match-partial}}]
{
  \Sigma; \Gamma \vdash e \synth \sigma \\
  \Sigma; \Gamma; \zeta \vdash \sigma, \shapes{\Sigma}(\sigma, \varnothing), \seq{b} : S, \tau, L \\
  \tau \neq \tyNever
}
{
  \Sigma; \Gamma; \zeta \vdash \exMatch{e}{\seq{b}} : S \merge_{\added{\Sigma}} \Assigns{\varnothing}
}

\end{mathpar}
The remaining type, the residual, reachability and exhaustiveness are the subject of
\secref{well-formedness-patterns}.

The \ruleName{defs} rule for a mutual def region of function definitions synthesises a type scheme for each definition
(\specfigref{typing-def}), binds all the function names of the mutual def region at the synthesised type schemes, and checks each
definition against its type scheme:
\begin{mathpar}
\small
\inferrule*[lab={\ruleName{defs}}]
{
  \Sigma; \Gamma \vdash d_i \synth \pi_i \quad (\forall i) \\
  \Delta' = \set{\bind{\declaresS(\seq{d})}{\seq{\pi}}} \\
  \Sigma; \Gamma \override \Delta' \vdash d_i \checks \pi_i \quad (\forall i)
}
{
  \Sigma; \Gamma; \zeta \vdash \seq{d} : \Assigns{\Delta'}
}

\end{mathpar}
For a definition $\exDef{f[\seq{\alpha}]}{\bind{\seq{x}}{\seq{\psi}}}{\exRetAnnot{\psi}}{s}$,
the auxiliary judgements $\Sigma; \Gamma \vdash d \synth \pi$ and $\Sigma; \Gamma \vdash d \checks \pi$ synthesise the type
scheme from the annotations and check the body $s$. Both rules extend the
context with $\TypeVar$ entries for the type parameters $\seq{\alpha}$. A type scheme without type parameters is identical to its \pyinline{Callable} type, so the \ruleName{defs} rule binds the name of a
function without type parameters as a variable of that type. The type schemes of functions with type parameters are instantiated
at a call, by the argument types (\ruleName{callee-scheme}), and when checked against a \pyinline{Callable} type, by that
type. Since annotations denote types, not type schemes, declarations and
assignments never bind variables at type schemes, only at types.

The body $s$ is checked with the
return type denoted by $\psi$, and must have static outcome $\Returns$ unless that type is equivalent to \pyinline{None},
so a function that declares a result definitely returns one. A return type such as $\tyUnion{\tyInt}{\tyNone}$
does not exempt the body from returning, nor permit a bare \pyinline{return}, for compatibility with mypy (\secref{python-compatibility}).

\subsection{Pattern matching}
\label{sec:exhaustiveness}
\label{sec:well-formedness-patterns}

To check a \pyinline{match} with scrutinee type $\sigma$, the auxiliary judgement
$\Sigma; \Gamma; \zeta \vdash \tau, L, \seq{b} : S, \tau', L'$ (\specfigref{typing-cases}) threads a \emph{remaining type}
$\tau$ and a \emph{residual} $L$ through the cases, both over-approximations of the values of $\sigma$ which are not
matched by any of the earlier cases. The residual is an ordered set of \emph{shapes} (\figref{shapes}), each denoting a
set of values of a single type, and starts from the shapes of $\sigma$; the remaining type is coarser, a type rather
than a set of shapes, and starts from $\sigma$ itself. Two auxiliary functions narrow a type $\tau$ by a pattern $p$: $\matchedAs{\Sigma}(\Gamma, \tau, p)$ is
the type to which $\tau$ is narrowed if $p$ matches, and $\remaining{\Sigma}(\Gamma, \tau, p)$ is the type remaining if
$p$ fails. Each case's pattern is checked against the remaining type, binding its pattern variables, and matched
against the residual. Reachability is decided by the
residual: a case is \emph{reachable} if its pattern matches at least one shape of the residual, and a match with an
unreachable case is ill-formed. For a reachable case, the body is checked under the bindings and its outcome merged
into $S$; the shapes left after matching and the type remaining if the pattern fails become the residual and remaining
type for the case that follows:
\begin{mathpar}
\small
\inferrule*[lab={\ruleName{match-case}}]
{
  \Sigma; \Gamma \vdash p \checks \tau \binding \Delta \\
  \Sigma; \Gamma \vdash L, p \matches K, L' \\
  \Sigma; \Gamma \override \Delta; \zeta \vdash s : R
}
{
  \Sigma; \Gamma; \zeta \vdash \tau, L, \exCase{p}{s} : \Assigns{\Delta} \override R, \remaining{\Sigma}(\Gamma, \tau, p), L'
}
\and
\inferrule*[lab={\ruleName{match-cases}}]
{
  \Sigma; \Gamma; \zeta \vdash \tau, L, b : S, \sigma, L' \\
  \Sigma; \Gamma; \zeta \vdash \sigma, L', \seq{b} : S', \tau', L^\dagger
}
{
  \Sigma; \Gamma; \zeta \vdash \tau, L, b \cons \seq{b} : S \merge_{\added{\Sigma}} S', \tau', L^\dagger
}

\end{mathpar}

A match is \emph{exhaustive} when its cases between them match every value of $\sigma$. Exhaustiveness is decided by
the final remaining type, which the \ruleName{match} rule requires to be $\tyNever$, with \ruleName{match-partial}
applying otherwise. The remaining type thus determines exhaustiveness as well as the types of pattern variables. The Python typing
specification leaves narrowing largely unspecified,\footnote{\url{https://typing.python.org/en/latest/spec/narrowing.html}}
so the remaining type emulates the narrowing of a match scrutinee observed in mypy
(\secref{python-compatibility}), which subtracts a class only for a
constructor pattern that matches every value of the class, and a tuple only through one refutable component, i.e.~one whose pattern can fail to match a value of the component's type.

\paragraph{Pattern checking.}
Checking a pattern against the remaining type (\ruleName{match-case}) rejects a tuple pattern at a type
with list values and a list pattern at a type with tuple values, and either at a type with range values, because matching either form of pattern
against the other kind of sequence or a range aborts with $\errExcluded$ (\secref{eval-patterns}). It also rejects a constructor
pattern at a type which does not determine the type arguments omitted by the pattern.

\paragraph{Shapes.}
\begin{figure}
\small
\added{\begin{tabular}[t]{rcll}
\grammarGroup{Heads} \\
$h$ & $::=$ & $\ell$ & (Literal) \\
    & $\mid$ & $C$ & (Class) \\
    & $\mid$ & $\headList{n}$ & (List of length $n$) \\
$H$ & $::=$ & $\set{h_1, \ldots, h_n}$ & (Excluded heads)
\end{tabular}
\hfill
\begin{tabular}[t]{rcll}
\grammarGroup{Shapes} \\
$k$ & $::=$ & $\shapeAny{\tau}{H}$ & (Rest) \\
    & $\mid$ & $\shapeAny{\patConstr{C[\seq{\tau}]}{k_1, \ldots, k_n}}{H}$ & (Constr) \\
    & $\mid$ & $\patTuple{k_1, \ldots, k_n}$ & (Tuple) \\
    & $\mid$ & $\shapeList{\tau}{k_1, \ldots, k_n}$ & (List) \\
    & $\mid$ & $\shapeDict{\tau}{\beta}{H}$ & (Dict) \\
$K$, $L$ & $::=$ & $\set{k_1, \ldots, k_n}$ & (Shape set)
\end{tabular}
}
\caption{\added{Heads and shapes}}
\label{fig:shapes}
\end{figure}
%
A shape (\figref{shapes}) denotes a set of values of a single type. A head is a literal, a class or a list
length, and denotes the values with that outermost form. The rest shape $\shapeAny{\tau}{H}$ denotes the
values of type $\tau$ whose head is not in the excluded set $H$. The other forms are constructor, tuple, list
and dictionary shapes, with shapes in place of sub-patterns; a constructor shape also has an excluded set,
containing the subclasses already matched. The auxiliary function $\shapes{\Sigma}(\tau, H)$ gives the
shapes of $\tau$ once the heads in $H$ are excluded.

\paragraph{Matching.}
The judgement $\Sigma; \Gamma \vdash L, p \matches K, L'$ matches a pattern $p$ against an ordered set of
shapes $L$, giving matched shapes $K$ and residual $L'$, and the judgement $\Sigma; \Gamma \vdash k, p
\matches K, L$ matches a pattern against a single shape $k$:
\begin{mathpar}
\small
\inferrule*[lab={\ruleName{pat-shapes}}]
{
  L' = \set{k \in L \mid \Sigma; \Gamma \vdash k, p \matches K_k, L_k} \\
  L' \neq \varnothing
}
{
  \Sigma; \Gamma \vdash L, p \matches \textstyle\bigdisjunion_{k \in L'} K_k, \; \textstyle\bigdisjunion_{k \in L'} L_k \disjunion (L \setminus L')
}

\and
\inferrule*[lab={\ruleName{pat-split}}]
{
  \Sigma; \Gamma \vdash k, p \splits K, L \\
  \Sigma; \Gamma \vdash K, p \matches K', L'
}
{
  \Sigma; \Gamma \vdash k, p \matches K', L \disjunion L'
}

\\
\inferrule*[lab={\ruleName{split-class}}]
{
  \Gamma \vdash q : \Type{C} \\
  \ancestors{\Sigma}(C) \cap H = \varnothing \\
  \Sigma \vdash C[\seq{\tau}'] \instanceAt \tau \\
  \meetTy{\Sigma}{\tau}{C[\seq{\tau}']} = D[\seq{\sigma}] \\
  \tyInst{\fields{\Sigma}(D)}{\seq{\sigma}} = \bind{\seq{x}}{\seq{\sigma}'} \\
  H' = \set{h \in H \mid \Sigma \vdash h : D[\seq{\sigma}]}
}
{
  \Sigma; \Gamma \vdash \shapeAny{\tau}{H}, \patConstr{q}{a} \splits \set{\shapeAny{\patConstr{D[\seq{\sigma}]}{\seq{k}}}{H'} \mid \seq{k} \in \shapes{\Sigma}(\seq{\sigma}')}, \shapes{\Sigma}(\tau, H \cup \set{C})
}

\\
\inferrule*[lab={\ruleName{pat-constr}}]
{
  \Gamma \vdash q : \Type{D} \\
  D \neq \clsRange \\
  D \in \ancestors{\Sigma}(C) \\
  \fields{\Sigma}(D) = \scheme{\seq{\alpha}}{\bind{\seq{x}}{\seq{\sigma}}} \\
  \fieldMap{\Sigma}(D, a) = \set{\bind{\seq{x}}{\seq{p}}} \\
  p'_i = \patWild \quad (\forall i) \\
  \Sigma; \Gamma \vdash \seq{k}, \seq{p} \mathbin{\concat} \seq{p}\,' \matches K, L
}
{
  \Sigma; \Gamma \vdash \shapeAny{\patConstr{C[\seq{\tau}]}{\seq{k}}}{H}, \patConstr{q}{a} \matches \set{\shapeAny{\patConstr{C[\seq{\tau}]}{\seq{k}\,'}}{H} \mid \seq{k}\,' \in K}, \set{\shapeAny{\patConstr{C[\seq{\tau}]}{\seq{l}}}{H} \mid \seq{l} \in L}
}

\end{mathpar}
The \ruleName{pat-shapes} rule matches the pattern against each shape of $L$. A case is reachable only if
its pattern matches at least one of the shapes, and the shapes which are not matched by the pattern pass into
the residual unchanged. A pattern matches a shape directly if the shape has the head tested by the pattern.
Otherwise \ruleName{pat-split} first splits the shape, by the judgement $\Sigma; \Gamma \vdash k, p \splits
K, L$, into shapes $K$ with that head and a residual $L$ with the head excluded. Splitting a rest shape at a
class $C$ (\ruleName{split-class}) gives constructor shapes for the instance of $C$ at the type of the shape,
with the shapes of the field types as components, and excludes $C$ from the residual; the instance
judgement, defined in \specfigref{instance}, recovers the type arguments omitted by the pattern. Shapes are split
only when a pattern requires it, because a class may refer to itself through its fields, and expanding every
shape in advance would not terminate. The \ruleName{pat-constr} rule matches a pattern for a class $D$
against a constructor shape for a subclass $C$ of $D$, matching the sub-patterns, with wildcards for omitted
fields, against the field shapes position by position.

The residual and the remaining type can differ. In the following program, the two cases between them match
every value of \pyinline{C}:

\begin{listingrow*}
\begin{listingcol}{0.45\textwidth}
\begin{python}[numbers=left, xleftmargin=1.5em]
@dataclass
class C:
    p: tuple[int, int] | None
\end{python}
\end{listingcol}
\begin{listingcol}{0.49\textwidth}
\begin{python}[numbers=left, firstnumber=4, xleftmargin=1.5em]
# error: f does not always return
def f(c: C) -> int:
    match c:
        case C(None):
            return 0
        case C((a, b)):
            return a + b
\end{python}
\end{listingcol}
\end{listingrow*}

\noindent The residual starts as the rest shape of \pyinline{C}. The first case splits it into a shape of
\pyinline{C} with a tuple field and a shape of \pyinline{C} with a \pyinline{None} field, and matches the
shape with the \pyinline{None} field. The second case matches the shape with the tuple field, leaving the
empty residual, so a third case would be unreachable.
Neither case matches every value of \pyinline{C}, however, so the remaining type is still \pyinline{C},
\ruleName{match-partial} applies, and the function is rejected as not always returning, following mypy. If
the second case were \pyinline{C(q)} instead, the remaining type after it would be $\tyNever$, but
\pyinline{q} would be bound at $\tyUnion{\tyTuple{\tyInt, \tyInt}}{\tyNone}$ rather than at $\tyTuple{\tyInt, \tyInt}$,
because pattern checking uses the remaining type.



The cases needed to exhaust a match therefore depend on the type of the scrutinee: a union is exhausted
by covering each disjunct, a class type by a pattern for the class itself, \pyinline{bool} by
\pyinline{True} and \pyinline{False}, and a dictionary type by the empty dictionary pattern, while \pyinline{int}, \pyinline{float}, \pyinline{str} and lists have
unboundedly many heads and need a variable or wildcard. A literal-typed scrutinee is exhausted by covering
its literals, so \replaced{$\exAssign{\pyinline{x}\exAnnot{\tyLiteral{\pyinline{42}}}}{\pyinline{42}}$}{\pyinline{x: Literal[42] = 42}} admits a match with a case for \pyinline{42} alone, while
\replaced{$\exAssign{\pyinline{x}\exAnnot{\tyInt}}{\pyinline{42}}$}{\pyinline{x: int = 42}} keeps the type open. Over
$\tyTuple{\tyBool, \tyBool}$ the cases $\patTuple{\pyinline{True}, \patWild}$,
$\patTuple{\patWild, \pyinline{False}}$ and $\patTuple{\pyinline{False}, \pyinline{True}}$ leave the empty residual,
although each covers only part of the type.

\subsection{Modules and programs}
\label{sec:well-formedness-modules}

\emph{Module checking} types a module's body and computes its signature $\Gamma$, the judgement
$\Sigma \vdash_{\mathsf{mod}} q : \Gamma \binding \Sigma'$ (\specfigref{typing-module}); it mirrors loading at run time
(\secref{operational-semantics}). The \ruleName{module} rule types the body as a top-level statement.
Module checking is defined inductively, so a program with an import cycle is ill-formed. Program
well-formedness (\specfigref{typing-module}) requires the main module to be well typed and hence,
recursively, every module that the program can load; a module that is never imported is not checked.
The rule \ruleName{module} types a module body under the context of $\mathcal{P}(\pyinline{builtins})$,
overridden by the imported context and then, as in \ruleName{def-body}, by an unbound entry for each name
declared in the body or bound by a pattern in it:
\begin{mathpar}
\small
\inferrule*[lab={\ruleName{module}}]
{
  \mathcal{M}(q) = \seq{\iota}\;t \\
  \mathcal{P}(\pyinline{builtins}) = (\Gamma, \rho) \\
  \Sigma; q \vdash \seq{\iota} : \Gamma' \binding \Sigma^\dagger \\
  \Sigma^\dagger; q; \Gamma \override \Gamma' \mathbin{\override} \scopeS(\declaresS(\seq{\iota}), t) \vdash t : \Assigns{\Delta} \binding \Sigma' \\
  (\dom(\Gamma') \cup \dom(\Delta)) \cap \dom(\submods(q)) = \varnothing
}
{
  \Sigma \vdash_{\mathsf{mod}} q : \submods(q) \override \Delta \binding \Sigma'
}

\end{mathpar}
The members of \pyinline{builtins} are thus in scope unqualified, and a name bound by the body
shadows a builtin throughout the module, so a use of the name before its declaration is rejected. Under mypy such a
use instead resolves to the builtin, following Python's run-time lookup. As with any scope, the names declared in the
module, including its imported names, must be distinct from one another and from any pattern variables in its
body. Since every module declares \pyinline{__name__} (\secref{modules-programs}), a module is ill-formed if its body
declares that name again, imports it, or binds it in a pattern.

The signature $\Gamma$ produced by module checking (\specfigref{typing-module}) comprises the classes, type aliases
and variables defined by a module, together with its intrinsic \pyinline{__name__} and a stub for each submodule.

